%% Courbes représentatives des fonctions cosinus et sinus sur [-4 ; 4] (manuel p. 101-102).
%% Cosinus : rouge, trait plein, fonction paire (courbe symétrique par rapport à l'axe des
%% ordonnées, rappelé par le trait pointillé rouge vertical).
%% Sinus : bleu, tirets, fonction impaire (courbe symétrique par rapport à O, marqué par un point).
\documentclass[tikz,border=4pt]{standalone}
\usepackage{liaisons}
\begin{document}
\begin{tikzpicture}[x=0.85cm,y=0.9cm,
  axe/.style={-{Stealth[length=5pt]}, line width=0.6pt},
  ccos/.style={solideA, line width=1.3pt, line join=round},
  csin/.style={solideB, line width=1.3pt, line join=round, dash pattern=on 4pt off 2.5pt},
  grad/.style={font=\scriptsize, fill=white, inner sep=1.2pt}]

%% ---------------- quadrillage ------------------------------------------
\draw[solideB!25, line width=0.3pt, xstep=0.5, ystep=0.5] (-4,-1.5) grid (4,1.5);

%% ---------------- axes et graduations ----------------------------------
\draw[axe] (-4.35,0) -- (4.4,0) node[anchor=north east, inner sep=3pt, font=\small] {$x$};
\draw[axe] (0,-1.75) -- (0,1.9);
\node[font=\small, anchor=north east, inner sep=2pt] at (-0.05,-0.05) {$O$};
\foreach \xv in {-4,-3,-2,-1,1,2,3,4} {
  \draw[line width=0.5pt] (\xv,0.08) -- (\xv,-0.08);
  \node[grad, anchor=north, inner sep=2pt] at (\xv,-0.09) {$\xv$};}
\foreach \yv in {-1,1} {
  \draw[line width=0.5pt] (0.08,\yv) -- (-0.08,\yv);
  \node[grad, anchor=east, inner sep=3pt] at (-0.09,\yv) {$\yv$};}

%% ---------------- abscisses remarquables (multiples de pi) --------------
\foreach \xv/\labr in {-3.1416/{-\pi}, -1.5708/{-\frac{\pi}{2}},
                       1.5708/{\frac{\pi}{2}}, 3.1416/{\pi}} {
  \draw[black!35, line width=0.5pt, dotted] (\xv,0) -- (\xv,-1.28);
  \node[font=\scriptsize, text=black!60, anchor=north, inner sep=2pt] at (\xv,-1.3) {$\labr$};}

%% ---------------- les deux courbes --------------------------------------
\draw[ccos] plot[domain=-4:4, samples=201, smooth] (\x,{cos(\x r)});
\draw[csin] plot[domain=-4:4, samples=201, smooth] (\x,{sin(\x r)});
\node[font=\small, text=solideA, anchor=south west, inner sep=2pt] at (0.18,1.02) {$\mathcal{C}_{\cos}$};
\node[font=\small, text=solideB, anchor=north east, inner sep=2pt] at (-1.75,-1.02) {$\mathcal{C}_{\sin}$};

%% ---------------- repères de symétrie -----------------------------------
\draw[solideA, line width=0.8pt, dotted] (0,-1.6) -- (0,1.75);   % axe de symétrie du cosinus
\fill[solideB] (0,0) circle (2pt);                               % centre de symétrie du sinus

%% ---------------- légende ------------------------------------------------
\begin{scope}[shift={(1.15,1.15)}]
  \draw[fill=white, draw=black!35, line width=0.5pt, rounded corners=2pt]
       (0,0) rectangle (2.75,0.72);
  \draw[ccos] (0.12,0.5) -- (0.6,0.5);
  \node[font=\scriptsize, anchor=west, inner sep=2pt] at (0.66,0.5) {$\mathcal{C}_{\cos}$ : paire};
  \draw[csin] (0.12,0.22) -- (0.6,0.22);
  \node[font=\scriptsize, anchor=west, inner sep=2pt] at (0.66,0.22) {$\mathcal{C}_{\sin}$ : impaire};
\end{scope}
\end{tikzpicture}
\end{document}
